\chapter{Implementation}
\label{Chapter5}
\lhead{Chapter 5. \emph{Implementation}}

CausalTrace is approximately 8{,}700 LOC; the kernel side is one 1{,}728-line
BCC file, the userspace pipeline is thirteen Python modules plus a loader
and an orchestrator. Table~\ref{tab:file-manifest} is the file manifest.

\begin{table}[h]
\centering
\caption[Source-tree file manifest]{File manifest by role. The single large
BCC file is deliberate: BCC~0.31's build model favours one translation unit
per program.}
\label{tab:file-manifest}
\footnotesize
\setlength\tabcolsep{3pt}
\begin{tabular}{@{}l r p{6.0cm}@{}}
\toprule
File & LOC & Role \\
\midrule
\texttt{kernel/causaltrace\_bcc.c}     & 1728 & Dispatcher + 6 handlers + Tier-2 probes + TC \\
\texttt{loader.py}                     &  467 & BCC compile, probe attach, TC pin/attach \\
\texttt{supervisor.py}                 &  209 & Daemon heartbeat watchdog \\
\texttt{tier3/daemon\_main.py}         &  411 & Detection-cycle orchestrator \\
\texttt{tier3/sheaf\_detector.py}      &  506 & Mahalanobis + sheaf + Rayleigh \\
\texttt{tier3/calibrate.py}            &  260 & PCA + whitener + CCA + thresholds \\
\texttt{tier3/enforcement\_engine.py}  &  665 & Graduated L0--L6 enforcement \\
\texttt{tier3/trust\_promoter.py}      &  133 & L4 stability promotion \\
\texttt{tier3/signal\_extractor.py}    &  130 & Bigram $\to$ 74-dim signal \\
\texttt{tier3/whitener, ema\_buffer,}\newline
\texttt{\phantom{tier3/}eigenmode\_analyzer}
                                       &  210 & Whitening, guarded EMA, eigenmode \\
\texttt{tier3/ringbuf\_monitor,}\newline
\texttt{\phantom{tier3/}verdict\_writer}
                                       &  237 & Ring-buffer reader, verdict writer \\
\texttt{run\_marathon\_evaluation.py}  & 1335 & Marathon orchestrator \\
\texttt{scripts/marathon\_analyze.py}  &  473 & Per-scenario analysis, Wilson CIs \\
\bottomrule
\end{tabular}
\end{table}

The modular \texttt{kernel/handler\_*.bpf.c} files in the repository are
reference implementations written during mid-review; they are not loaded at
runtime.

%─────────────────────────────────────────────────────────────────────────────
\section{The verifier fights}
\label{sec:verifier}
%─────────────────────────────────────────────────────────────────────────────

Linux~6.17's BPF verifier rejected the first compile of the dispatcher for
three separate reasons. Each bug had a terse error message, a reproducible
trigger, and a specific workaround that now lives inline in the kernel
source.

\paragraph{Bug 1: direct \texttt{pt\_regs} dereference past offset~104.}
The standard BCC idiom
\mbox{\texttt{struct pt\_regs *regs = (struct pt\_regs *)ctx->args[0]; long
arg = regs->di;}} is rejected on kernel~6.17 at offset~112 with
\texttt{invalid mem access 'inv'}. The verifier's register-state model
tracks memory accesses up to 104 bytes past the context; reads at 112 bytes
violate that bound. Every register read is routed through
\texttt{bpf\_probe\_read\_kernel(\&\_\_arg, sizeof(\_\_arg), \&regs->di)}
instead---one indirect-call overhead per argument in exchange for a program
that actually loads.

\paragraph{Bug 2: Clang fuses bounded zero-loops into \texttt{memset}.}
The Count-Min Sketch reset at the start of each detection cycle zeros all
$4 \times 128 = 512$ counters. A naive nested loop compiles to a
\texttt{memset} call on the 6.17 Clang toolchain; \texttt{memset} is not a
BPF helper and the program is rejected with \texttt{invalid function call}.
The fix is a \texttt{\#pragma unroll} block containing volatile stores:
\mbox{\texttt{*(volatile u32 *)\&sketch->counters[...]~=~0}}. The
\texttt{volatile} blocks memset fusion; the unroll keeps the loop visible
to the verifier as a sequence of stores.

\paragraph{Bug 3: unbounded \texttt{strncpy}-style loops.}
The basename extraction in \texttt{handle\_execve} scans up to 127 bytes of
the filename for the last \texttt{/}. A plain \texttt{for} loop with a
runtime bound fails the verifier's termination check. The bound is hoisted
to a compile-time constant and the loop is unrolled with
\texttt{\#pragma unroll}; we accept a code-size increase in exchange for
verifier acceptance. The same pattern applies to the path-prefix matching
in \texttt{handle\_file}.

\paragraph{Other constraints.} \texttt{strcmp} and \texttt{memcmp} are not
available; every comparison is byte-by-byte.
\texttt{bpf\_map\_lookup\_elem} returns a possibly-null pointer that the
verifier insists be checked on every access. \texttt{printk} output is
limited to 64 bytes; all diagnostics go through the
\texttt{alerts\_rb}~/~\texttt{telemetry\_rb} ring buffers.

%─────────────────────────────────────────────────────────────────────────────
\section{BPF maps and handler attachment}
\label{sec:maps}
%─────────────────────────────────────────────────────────────────────────────

Table~\ref{tab:bpf-maps} lists the maps that matter for the thesis
discussion; Appendix~\ref{AppendixB} carries the full schema.
Table~\ref{tab:handler-attach} maps each handler to its tracepoint or
kprobe.

\begin{table}[h]
\centering
\caption[Key BPF maps]{Key BPF maps by role. 18 maps total; the
appendix reproduces the full schema with key/value sizes.}
\label{tab:bpf-maps}
\footnotesize
\setlength\tabcolsep{4pt}
\begin{tabular}{@{}l l r p{5.9cm}@{}}
\toprule
Map & Type & Size & Role \\
\midrule
\texttt{prog\_array}         & \texttt{PROG\_ARRAY} & 512    & Tail-call dispatch table by syscall nr \\
\texttt{bigram\_sketch\_map} & \texttt{HASH}        & 256    & Per-cgroup CMS ($4 \times 128$) \\
\texttt{container\_behavior} & \texttt{HASH}        & 256    & Per-cgroup behaviour bits + $\mathrm{bit\_ts}[8]$ \\
\texttt{rate\_map}           & \texttt{HASH}        & 256    & Per-cgroup fork rate + 2nd-derivative state \\
\texttt{client\_trust}       & \texttt{HASH}        & 4096   & Per-IP trust state \\
\texttt{tid\_client\_ip}     & \texttt{LRU\_HASH}   & 16384  & Per-TID attribution to remote client \\
\texttt{connection\_context} & \texttt{LRU\_HASH}   & 8192   & Per-socket byte accumulator \\
\texttt{drop\_ip\_list}      & \texttt{HASH}        & 4096   & Burn-list, 5-min TTL per IP \\
\texttt{alerts\_rb}          & \texttt{RINGBUF}     & 64\,K  & Critical alerts to userspace \\
\texttt{telemetry\_rb}       & \texttt{RINGBUF}     & 256\,K & Connection / exec events \\
\texttt{verdict\_map}        & \texttt{HASH}        & 256    & Tier-3 KILL decrees \\
\texttt{enforce\_level\_map} & \texttt{HASH}        & 256    & Graduated L0--L6 level + TTL \\
\bottomrule
\end{tabular}
\end{table}

\begin{table}[h]
\centering
\caption[Handler attachment]{Tier-1 handlers dispatch via the tail-call
table; the Tier-2 probes and TC classifiers attach directly to kernel
functions rather than syscall entry, so the ``syscalls'' column lists
the kernel function they hook instead.}
\label{tab:handler-attach}
\small
\setlength\tabcolsep{4pt}
\begin{tabular}{@{}l l l@{}}
\toprule
Handler & Tracepoint / kprobe & Syscall~nr~/~kernel~fn \\
\midrule
\texttt{dispatcher}             & \texttt{raw\_tracepoint/sys\_enter}   & all syscalls \\
\texttt{handle\_dup2}           & tail-call                             & 33, 292 \\
\texttt{handle\_fork}           & tail-call                             & 56, 435 \\
\texttt{handle\_execve}         & tail-call                             & 59 \\
\texttt{handle\_file}           & tail-call                             & 257 \\
\texttt{handle\_file\_open}     & tail-call                             & 2 \\
\texttt{handle\_privesc}        & tail-call                             & 101, 105, 272, 308 \\
\texttt{trace\_connect\_return} & \texttt{kretprobe/tcp\_v4\_connect}    & \texttt{tcp\_v4\_connect} \\
\texttt{trace\_accept\_return}  & \texttt{kretprobe/inet\_csk\_accept}   & \texttt{inet\_csk\_accept} \\
\texttt{trace\_sendmsg}         & \texttt{kprobe/tcp\_sendmsg}            & \texttt{tcp\_sendmsg} \\
\texttt{trace\_recvmsg}         & \texttt{kprobe/tcp\_recvmsg}            & \texttt{tcp\_recvmsg} \\
\texttt{tc\_drop\_ingress}      & \texttt{SCHED\_CLS} veth ingress       & packet hook \\
\texttt{tc\_drop\_egress}       & \texttt{SCHED\_CLS} veth egress        & packet hook \\
\bottomrule
\end{tabular}
\end{table}

The loader (\texttt{loader.py}) compiles the BCC source, installs the six
handlers into \texttt{prog\_array}, attaches the dispatcher and the Tier-2
kprobes, pins the two TC classifier programs under
\texttt{/sys/fs/bpf/causaltrace}, and attaches them to every monitored
container's host-side veth with \texttt{tc filter add dev <veth> ingress
bpf direct-action}. Signal handlers (\texttt{SIGTERM}, \texttt{SIGINT}) and
an \texttt{atexit} callback detach every probe and TC filter cleanly; the
\texttt{--cleanup} mode reports any stale programs from a crashed earlier
run.

%─────────────────────────────────────────────────────────────────────────────
\section{Tier-3 daemon and graduated enforcement}
\label{sec:daemon}
%─────────────────────────────────────────────────────────────────────────────

\texttt{tier3/daemon\_main.py} is the runtime entry point. Each detection
cycle: read the Count-Min Sketches, snapshot the behaviour bits, drain the
telemetry ring buffer, whiten and project through every edge's restriction
maps, compute Mahalanobis energies and the Rayleigh quotient, run compound
confirmation, and write a verdict record. The cycle cost is small relative
to the 5-second budget.

The enforcement engine (\texttt{tier3/enforcement\_engine.py}) maps each
verdict to one of seven levels:
\textbf{L0~OBSERVE}~(telemetry only),
\textbf{L1~DENY}~(specific \texttt{connect}/\texttt{openat}/\texttt{execve}
calls fail via \texttt{bpf\_override\_return}),
\textbf{L2~SEVER}~(targeted flow drops via \texttt{drop\_ip\_list}),
\textbf{L3~THROTTLE}~(rate-limit to specific destinations),
\textbf{L4~FIREWALL}~(allow-list only calibrated destinations),
\textbf{L5~DRAIN}~(reduce egress rate cap),
\textbf{L6~QUARANTINE}~(block all outbound connects with
\texttt{ENETUNREACH}).
Every \texttt{enforce\_level\_map} entry carries a TTL so enforcement
self-expires if the threat has passed; the engine refreshes TTLs on
continued anomaly.

%─────────────────────────────────────────────────────────────────────────────
\section{Marathon orchestrator}
\label{sec:marathon}
%─────────────────────────────────────────────────────────────────────────────

\texttt{run\_marathon\_evaluation.py} sequences the evaluation in four
phases. Phase~1 runs dense calibration traffic through wrk plus
inter-service docker-exec patterns. Phase~2 replays a fixed, seed-42
permutation of 150 in-distribution attacks plus 5 held-out S11 fileless
\texttt{memfd\_create} injections against CausalTrace. Phase~3 tears down
CausalTrace, starts Falco (stock rules, then tuned rules), and replays the
identical 155-injection sequence. Phase~4 does the same for Tetragon
(stock policies, then CTEval \texttt{TracingPolicies}). All three detectors
see the identical attack order under the identical concurrent background
load, so baseline comparison is not confounded by hardware or traffic
variability. The analyser \texttt{scripts/marathon\_analyze.py} reads the
per-detector JSONL outputs and produces the per-scenario detection matrix,
Wilson 95\,\% confidence intervals, time-to-kill percentiles, and the
baseline-comparison CSV that Chapter~\ref{Chapter6} consumes.

Attack scripts live in \texttt{attacks/} (eleven scenarios plus two evasion
variants); Appendix~\ref{AppendixA} tabulates the full catalogue with
MITRE~ATT\&CK annotations.
